%---------------------------------------------------------------------
\subsection{Biomechanical and Physiological Foundations}

In a full-crimp grip (DIP joint hyperextended), biomechanical modelling estimates A2 pulley forces at approximately $2.7\times$ the external fingertip force (${\sim}255$~N at ${\sim}96$~N fingertip load) \cite{VQLVM06}, while A2 loading in crimp is approximately 36 times greater than in the open-hand (slope) position---a grip-type comparison, not a load multiplier. Cadaveric studies confirm high A2 pulley loads in crimp (380--700~N across finger positions) \cite{Sch01}, accounting for the predominance of A2 ruptures in climbing. Half-crimp and open-hand grips generate substantially lower pulley loads, underpinning standard advice against routine full-crimp hangboarding. Middle and ring fingers bear disproportionately higher loads than index and little fingers \cite{VQPC08}$^{\dagger}$, consistent with clinical injury patterns.

Climbing physiology reviews \cite{Wat04, MSB+07, FSS+16} establish three points central to hangboard training design: MFS is the single strongest measurable predictor of redpoint grade in trained cohorts; forearm endurance and oxidative capacity independently predict sustained-route performance; and these qualities are specific enough that general upper-body resistance training does not substitute for direct finger loading.

\smallskip
\noindent$^{\dagger}$\textit{Vigouroux et al. \cite{VQPC08} is a conference abstract; the force-distribution finding is consistent with the broader biomechanical literature but should be treated as supporting rather than primary evidence.}
